\documentclass[11pt,a4paper]{article}
\usepackage[margin=1in]{geometry}
\usepackage[T1]{fontenc}
\usepackage{lmodern}
\usepackage{microtype}
\usepackage{hyperref}
\usepackage{array}
\usepackage{longtable}
\hypersetup{
  colorlinks=true,
  linkcolor=black,
  urlcolor=black,
  pdftitle={OC Core 1.3.2 Methods and Reproducibility Companion},
  pdfauthor={Alexander Yashin},
  pdfsubject={Ontology of Continua Core 1.3.2 methods and reproducibility companion},
  pdfkeywords={Ontology of Continua, OC Core 1.3.2, reproducibility, methods}
}
\title{\textbf{OC Core 1.3.2 Methods and Reproducibility Companion}\\
\large Release-Governed Evidence, Replay, and Audit Route}
\author{Alexander Yashin\\Independent Researcher, Leipzig/Halle, Germany\\ORCID 0009-0008-6166-0914}
\date{Version v1.3.2; release date: 28.04.2026}
\begin{document}
\maketitle

\section*{Scope}
This companion describes how OC Core 1.3.2 connects scientific claims to reproducible release artifacts. It is a release-bound methods surface, not a separate theory authority.

\section*{Methodological Rule}
Every release-visible scientific assertion must be mapped to one of the following support routes:
\begin{itemize}
  \item proof or source-audited theorem support;
  \item replay pass with fixed input, code route, output hash, and tolerance envelope;
  \item bounded operational support with explicit falsifier and claim ceiling;
  \item strict salvage state when a line is historically relevant but not Core authority;
  \item owner-gated publication or value externalization state when the blocker is not a scientific theorem.
\end{itemize}

\section*{Reproducibility Surfaces}
\begin{longtable}{>{\raggedright\arraybackslash}p{0.28\linewidth}>{\raggedright\arraybackslash}p{0.64\linewidth}}
\textbf{Surface} & \textbf{Reproducibility role}\\
\hline
Release machine package/evaluate & Rebuilds package metadata, manifests, checksums, ZIP integrity report, and no-send verdict.\\
PDF quality gate & Extracts text from primary PDFs and fails on unresolved references, local paths, stale version strings, stale closure language, missing release identity, or missing master dedication.\\
K0--K12 projection state & Records whether each domain projection is replay-backed, source-audited, or bounded as structural/operator signature.\\
DRT strict salvage & Keeps DRT outside Core authority unless workflow, figures, PDF, hashes, and Core-gate reconciliation pass.\\
Tsukrov repair & Preserves private reviewer-source material outside the package while exposing public-safe response routes and release references.\\
Parfitian Cerberus & Blocks release readiness when ethical/release-governance findings exceed the configured threshold.\\
LRGEF & Controls publication locks, owner approval, provenance, source freshness, and release artifact inventory.\\
\end{longtable}

\section*{Prediction and Quantitative Language}
The word ``prediction'' is release-permitted only when the corresponding artifact includes protocol, code, data or fixture, output hash, tolerance envelope, and falsifier. Without those components, OC Core 1.3.2 uses structural or operator-signature language.

\section*{Package Safety}
The public package excludes private paths, credentials, local machine references, private reviewer-source material, private model transcripts, placeholder source templates, and unsupported publication actions. The release ZIP is built from an explicit bundle manifest and a placeholder-payload gate.

\section*{Expected Local Verification}
Before owner review, the release must pass:
\begin{itemize}
  \item release machine package in no-publish mode;
  \item release machine dry-run evaluation;
  \item Parfitian Cerberus tests;
  \item PDF text extraction scans;
  \item public-safety leak scans;
  \item version consistency across visible pages, metadata, manifests, and release dossier.
\end{itemize}

\end{document}
